% Part: many-valued-logic
% Chapter: infinite-valued-logics
% Section: lukasiewicz

\documentclass[../../../include/open-logic-section]{subfiles}

\begin{document}

\olfileid[fr]{mvl}{inf}{luk}

\olsection{La logique de \L ukasiewicz}

La matrice de \L ukasiewicz conserve le minimum et le maximum
pour la conjonction et la disjonction. Le conditionnel dépend
de l'écart entre les valeurs de son antécédent et de son conséquent.

\begin{defn}\ollabel{def:lukasiewicz}
La logique de \L ukasiewicz à une
infinité de valeurs~$\LogLuk[\infty]$
est définie par la matrice suivante :
\begin{enumerate}
  \item Le langage propositionnel usuel
  $\Lang L_0$ avec $\lnot$, $\land$,
  $\lor$ et $\lif$.
  \item L'ensemble des valeurs de vérité
  $V_\infty$.
  \item $1$ est la seule valeur désignée,
  c'est-à-dire $V^+=\{1\}$.
  \item Les fonctions de vérité sont :
  \begin{align*}
    \tf{\lnot}[\LogLuk](x) &=1-x\\
    \tf{\land}[\LogLuk](x,y) &=\min(x,y)\\
    \tf{\lor}[\LogLuk](x,y) &=\max(x,y)\\
    \tf{\lif}[\LogLuk](x,y)
      &=\min(1,1-(x-y))
        =\begin{cases}
          1 & \text{si }x\le y,\\
          1-(x-y) & \text{sinon.}
        \end{cases}
  \end{align*}
\end{enumerate}
La logique de \L ukasiewicz à $m$
valeurs est définie de la même façon,
à ceci près que $V=V_m$.
\end{defn}

\begin{prop}
  La logique $\LogLuk[3]$ définie
  par la \olref[thr][luk]{def:lukasiewicz}
  coïncide, sous l'identification
  $\True=1$, $\Undef=1/2$ et
  $\False=0$, avec la logique
  $\LogLuk[3]$ définie par la
  \olref{def:lukasiewicz}.
\end{prop}

\begin{proof}
  Il suffit de comparer les tables
  de vérité des connecteurs données
  dans la \olref[thr][luk]{def:lukasiewicz}
  avec celles que déterminent les
  équations de la
  \olref{def:lukasiewicz} :
  \begin{center}
    \begin{tabular}{c|c}
      $\tf{\lnot}$ & \\
      \hline
      $1$ & $0$ \\
      $1/2$ & $1/2$ \\
      $0$ & $1$
    \end{tabular}
    \quad
    \begin{tabular}{c|ccc}
      $\tf{\land}[\LogLuk[3]]$ & $1$ & $1/2$ & $0$ \\
      \hline
      $1$ & $1$ & $1/2$ & $0$ \\
      $1/2$ & $1/2$ & $1/2$ & $0$ \\
      $0$ & $0$ & $0$ & $0$
    \end{tabular}
    \\[2ex]
    \begin{tabular}{c|ccc}
      $\tf{\lor}[\LogLuk[3]]$ & $1$ & $1/2$ & $0$ \\
      \hline
      $1$ & $1$ & $1$ & $1$ \\
      $1/2$ & $1$ & $1/2$ & $1/2$ \\
      $0$ & $1$ & $1/2$ & $0$
    \end{tabular}
    \quad
    \begin{tabular}{c|ccc}
      $\tf{\lif}[\LogLuk[3]]$ & $1$ & $1/2$ & $0$ \\
      \hline
      $1$ & $1$ & $1/2$ & $0$ \\
      $1/2$ & $1$ & $1$ & $1/2$ \\
      $0$ & $1$ & $1$ & $1$
    \end{tabular}
  \end{center}
\end{proof}

\begin{prop}\ollabel{prop:luk-infty-m}
  Si $\Gamma\Entails[\LogLuk[\infty]] !B$,
  alors
  $\Gamma\Entails[\LogLuk[m]] !B$
  pour tout~$m\ge2$.
\end{prop}

\begin{proof}
  Exercice.
\end{proof}

\begin{prob}
  Démontrer la
  \olref[mvl][inf][luk]{prop:luk-infty-m}.
\end{prob}

La réciproque vaut lorsque $\Gamma$ est fini. Dans un
contre-modèle rationnel, les valeurs des variables de $\Gamma$
et de $!B$ ont alors un dénominateur commun $d$. Les fonctions
de vérité se restreignent à $V_{d+1}$. En y conservant ces
valeurs et en donnant $0$ aux autres variables, on obtient
un contre-modèle dans cette matrice finie.

Pour des prémisses infinies, prenons des variables distinctes
$p_0,p_1,\ldots$ et écrivons $!A\liff !B$ pour
$(!A\lif !B)\land(!B\lif !A)$. Ce biconditionnel vaut $1$
exactement lorsque les deux valeurs coïncident. La prémisse
$p_0\liff\lnot p_0$ impose la valeur $1/2$ à $p_0$.
Pour chaque $i\geq 0$, ajoutons
\[
 p_i\liff(\lnot p_{i+1}\lif p_{i+1}).
\]
Le membre de droite a pour valeur le double de celle de
$p_{i+1}$, plafonné à $1$. Par récurrence, les prémisses
imposent donc la valeur $2^{-(i+1)}$ à $p_i$. Elles sont
satisfaites dans $V_\infty$ et ne le sont dans aucun $V_m$,
qui ne peut contenir toutes ces valeurs. Elles entraînent
ainsi $p_0$ dans chaque $V_m$, tandis que dans $V_\infty$
la valeur de $p_0$ est $1/2$.

La source présente la logique de
\L ukasiewicz à une infinité de valeurs
comme la logique floue la plus répandue.
Si l'on définit plutôt le conditionnel
par $\lnot !A\lor !B$, en conservant
les autres fonctions de vérité,
on obtient une logique forte de Kleene
à une infinité de valeurs.

\begin{prob}
  Montrer que
  $(p\lif q)\lor(q\lif p)$ est
  une tautologie de~$\LogLuk[\infty]$.
\end{prob}

\end{document}
